\documentclass[10pt]{article}
\usepackage[letterpaper,margin=0.9in]{geometry}
\usepackage{lmodern}
\usepackage[T1]{fontenc}
\usepackage[utf8]{inputenc}
\usepackage{microtype}
\usepackage{parskip}
\usepackage{enumitem}
\setlist{nosep,leftmargin=1.4em}
\usepackage{booktabs}
\usepackage{tabularx}
\usepackage{xcolor}
\definecolor{lrblue}{RGB}{22,80,145}
\usepackage{titlesec}
\titleformat{\section}{\large\bfseries\sffamily\color{lrblue}}{\thesection}{0.8em}{}
\titlespacing{\section}{0pt}{0.9em}{0.3em}
\usepackage{fancyhdr}
\pagestyle{fancy}
\fancyhf{}
\fancyhead[L]{\small\sffamily Lightrider Inc: Requests to wolfSSL}
\fancyhead[R]{\small\sffamily September 17, 2026 (rev.\ 3)}
\fancyfoot[C]{\small\thepage}
\newcommand{\code}[1]{\texttt{\small #1}}
\newcolumntype{L}{>{\raggedright\arraybackslash}X}

\begin{document}

{\LARGE\sffamily\bfseries Requests to wolfSSL}\\[0.2em]
{\sffamily Lightrider Inc, Veloce PQC SDK. Certificate \#4718, module
v5.2.1, bundle \code{wolfssl-5.9.2-commercial-fips-linuxv5.2.1},
Agreement v.01-24.}

\section{Our configuration, for your verification}

\begin{small}
\begin{tabularx}{\textwidth}{@{}lL@{}}
\toprule
Linux build & \code{./configure --enable-fips=v5}; make, fips-hash.sh,
  make, testwolfcrypt. No \code{--enable-wolfEntropy}, per your 2026-08-27
  guidance that wolfEntropy was never tested with v5.2.1. \\
Windows build & UG 2.2.45 (OE49): \code{wolfssl-fips.sln},
  \code{DLL Release | x64}, Visual Studio 2026, no PAA, your supplied
  \code{user\_settings.h} (\code{HAVE\_FIPS\_VERSION 5 / MINOR 2 / PATCH 1}). \\
Seeding & Custom generate-seed callback registered with
  \code{wc\_SetSeed\_Cb} (FIPS FAQ \S2). Revised 2026-09-17: the callback
  reads the processor's \code{RDSEED} instruction (conditioned output of the
  CPU's SP~800-90B noise source) and delivers it directly to the module's
  Hash\_DRBG, with SP~800-90B RCT (cutoff 4) and APT (13/512) health tests at
  $H = 8$ bits/sample and a 1024-sample startup test; fails closed on any
  error or when RDSEED is absent. Operating system DRBG output
  (\code{getrandom}, \code{BCryptGenRandom}) is no longer the default: it is
  an SP~800-90C RBGC chain and is available only by explicit configuration,
  reported with no security-strength claim. Documented under legacy
  IG~9.3.A; no ESV certificate claimed; caveat ``no assurance of the minimum
  strength of generated SSPs'' acknowledged. \\
Distribution & wolfCrypt as object code only inside Veloce installers;
  notices retained; source confined to a restricted private repository. \\
\bottomrule
\end{tabularx}
\end{small}

\section{Please verify or confirm}

\begin{enumerate}
  \item The Linux and Windows build configurations above are approved
        configurations for module v5.2.1 under the User Guide / COGM.
  \item The revised seeding design above (RDSEED direct to the module
        DRBG) raises no concern for v5.2.1 under legacy IG~9.3.A, and the
        earlier operating-system-DRBG seed source is correctly understood as
        an SP~800-90C RBGC construction that should not be the shipped
        default.
  \item Authenticode signing of the FIPS DLL or of the MSI installer does
        not interfere with the in-core HMAC-SHA-256 integrity check.
\end{enumerate}

\section{Please notify us}

\begin{enumerate}
  \item When \textbf{OE49} (Windows 11 Pro, Lenovo ThinkPad T14, Intel
        i7-1260P) and \textbf{OE36} (Ubuntu 24.04 on Docker 28.5 / Hyper-V,
        Xeon Platinum 8272CL, Azure Standard\_D2s\_v4; CAVP validation
        36918) appear on the published \#4718 certificate and security
        policy, with any expected CMVP timeframe. Until publication we ship
        Windows with validation status ``pending publication''.
\end{enumerate}

\section{Please advise on entropy source validation}

Your review noted that \code{getrandom} and \code{BCryptGenRandom} return the
output of the operating system's SP~800-90A DRBG, so seeding the module DRBG
from them forms an SP~800-90C RBGC construction whose randomness source is
not validated on our operational environments, and recommended seeding from
raw SP~800-90B entropy instead. We have moved the default seed source to
\code{RDSEED} accordingly. Under IG~9.3.A Additional Comment~11 an ESV
certificate counts only when its number appears on the module certificate
and in the security policy, so any ESV credit for Veloce depends on the
module side. Please advise on the following.

\begin{enumerate}
  \item \textbf{RDSEED as the external seed source.} For module v5.2.1
        under legacy IG~9.3.A, is a generate-seed callback that delivers
        \code{RDSEED} output directly to the Hash\_DRBG the seed design you
        recommend? Should the callback apply SP~800-90B RCT/APT health tests
        (as ours does), condition the output further (for example through a
        vetted conditioning function with oversampling), or deliver it
        unmodified? Which entropy estimate per 8-bit sample should the
        health-test cutoffs assume for Intel and AMD RDSEED?
  \item \textbf{RDRAND and TPM.} We read \code{RDRAND} and
        \code{TPM2\_GetRandom} as DRBG outputs, so both would again form
        RBGC chains, with a FIPS-validated TPM being a chain from a validated
        RBG. Please confirm or correct, and state whether a validated TPM is
        an acceptable alternative source where RDSEED is unavailable (arm64,
        some virtual machines).
  \item \textbf{Windows CNG.} On Windows~11, \code{BCryptGenRandom} is
        served by Microsoft's FIPS-validated Cryptographic Primitives
        Library. Is an RBGC construction from that validated RBG acceptable
        for the module on OE49, and what documentation would the security
        policy or user guide expect?
  \item \textbf{Entropy assessment path.} Whether wolfSSL can point to, or
        obtain under a commercial engagement, an ESV certificate or an
        SP~800-90B entropy assessment covering RDSEED on the Windows~11 x86-64
        (OE49) and Ubuntu x86-64 operational environments, and whether such a
        certificate would be cited in a wolfCrypt security policy so that
        IG~9.3.A Additional Comment~11 is satisfied.
  \item \textbf{ESV coverage of future modules.} The ESV certificate and
        operating-environment list that will back the v5.2.4 module and the
        FIPS v7 ``PQ-FS'' module, and whether Windows~11 x86-64 and
        Ubuntu~24.04 x86-64 will be covered. ESV~\#E335 currently lists one
        OE (AEDS Software Distribution 2.0, Xeon E3-1275 v6) with reuse
        restricted to vendor.
  \item \textbf{Third-party sources.} Whether wolfSSL will cite a
        third-party ESV certificate (a Lightrider-owned entropy source or an
        IDQ Quantis QRNG) in a wolfCrypt module security policy under a
        commercial engagement, or alternatively offer an OEM module
        validation for Lightrider.
  \item \textbf{Resubmission effect.} Whether the in-process resubmission of
        \#4718 for OE49 and OE36 changes the module's legacy IG~9.3.A status
        or its entropy caveat.
\end{enumerate}

\section{Please paper on the Order Form}

\begin{enumerate}
  \item Veloce as the named Program under \S3: a PQC SDK exposing
        cryptographic services through a local agent, not a linkable
        cryptography library.
  \item Commercial coverage for our PQC provider: standalone compilation of
        \code{wc\_mlkem.*}, \code{wc\_mldsa.*}, and \code{sha3.c} from the
        public \code{v5.9.2-stable} tree (the commercial bundle strips
        SHAKE), shipped as object code only.
  \item Budget and timing for the FIPS v7 ``PQ-FS'' upgrade (\S10 treats it
        as an upgrade requiring a new Order Form).
\end{enumerate}

\end{document}
